% !TEX TS-program = pdflatex
% !TeX program = pdflatex
% !TEX encoding = UTF-8
% !TEX spellcheck = en_US

\documentclass[xcolor=table]{beamer}

\usepackage{../extra/beamer/karimnlp}


%\usepackage{fullpage}
%\usepackage[left=2.8cm,right=2.2cm,top=2 cm,bottom=2 cm]{geometry}
\setbeamersize{text margin left=10pt,text margin right=10pt}



%\usepackage[T1]{fontenc}
\usepackage{textcomp}
\usepackage[utf8]{inputenc}
%\usepackage[french,english]{babel}
\usepackage{arabtex}
\usepackage{txfonts}
\usepackage{tipa}
\usepackage[]{graphicx}
\usepackage{multirow}
\usepackage{hyperref}
\usepackage{colortbl}
\usepackage{tabularray}
\usepackage{listingsutf8}
\usepackage{wrapfig}
\usepackage{multicol}
\usepackage[export]{adjustbox} %for images in table, also for frame
\usepackage[many]{tcolorbox}
\usepackage{wasysym}
\usepackage[lined]{algorithm2e}
\usepackage{alltt} %verbatim with commands
\usepackage{longtable}
\usepackage{tabu}

\usepackage{CJKutf8}

\usepackage{amsmath,amssymb, amsfonts} 

%\usepackage{newtxtext,newtxmath}


%\usepackage{acolor}


\definecolor{lightblue}{HTML}{D0D2FF}
\definecolor{lightyellow}{HTML}{FFFFAA}
\definecolor{darkblue}{HTML}{0000BB}
\definecolor{olivegreen}{HTML}{006600}
\definecolor{darkgreen}{HTML}{008B45} %009B55
\definecolor{violet}{HTML}{6600CC}
\definecolor{deeppink}{HTML}{FF1493}
\definecolor{orangey}{HTML}{FFBB00}


\makeatletter
\def\beamer@calltheme#1#2#3{%
	\def\beamer@themelist{#2}
	\@for\beamer@themename:=\beamer@themelist\do
	{\IfFileExists{\beamer@themelocation/#3\beamer@themename.sty}{%
	\usepackage[{#1}]{\beamer@themelocation/#3\beamer@themename}}{\usepackage[{#1}]{#3\beamer@themename}}
		
	}%
}

\def\usefolder#1{
	\def\beamer@themelocation{#1}
}
\def\beamer@themelocation{}
\makeatother

\usefolder{../../extra/beamer-themes}

\usetheme{Karim} % Antibes Boadilla Warsaw Karim

\beamertemplatenavigationsymbolsempty

\let\oldcite\cite
\renewcommand{\cite}[1]{{\bfseries\color{blue}\oldcite{#1}}}

\tcbuselibrary{listings}


%\renewcommand{\baselinestretch}{1.5}

\def\supit#1{\raisebox{0.8ex}{\small\it #1}\hspace{0.05em}}

\AtBeginSection{%
	\begin{frame}
		\sectionpage
	\end{frame}
}

\newcommand{\rottext}[2]{%
	\rotatebox{90}{%
	\begin{minipage}{#1}%
		\raggedleft#2%
	\end{minipage}%
	}%
}


\institute{ %
Laboratoire de la Communication dans les Systèmes Informatiques (LCSI)
\par
École  nationale Supérieure d'Informatique (ESI), Algiers, Algeria
}
\author[Abdelkrime Aries (2023/2024)] %
{Abdelkrime Aries}
%\titlegraphic{\includegraphics[height=1cm]{../img/esi-logo.png}%\hspace*{4.75cm}~


\date{Academic year: 2023/2024} %\today

\titlegraphic{%
	\includegraphics[height=1cm]{../../img/misc/esi-logo.png}%
	\hspace{2cm}%
	\includegraphics[height=1cm]{../../img/misc/lcsi-logo.png}%
	\hspace{2cm}%
	\includegraphics[height=1.5cm]{../../img/misc/esi.nlp-logo.png}%
}

%\setbeamertemplate{headline}{}

\newcommand{\kurl}[1]{{\scriptsize\bfseries\color{orangey}\url{#1}}}

\newcommand{\keyword}[1]{\textcolor{red}{\bfseries\itshape #1}}
\newcommand{\expword}[1]{\textcolor{olivegreen}{#1}}
\newcommand{\optword}[1]{\textcolor{violet}{\bfseries #1}}

\makeatletter
\newcommand\mysphere{%
	\parbox[t]{10pt}{\raisebox{0.2pt}{\beamer@usesphere{item projected}{bigsphere}}}}
\makeatother

%\let\oldtabular\tabular
%\let\endoldtabular\endtabular
%\renewenvironment{tabular}{\rowcolors{2}{white}{lightblue}\oldtabular\rowcolor{blue}}{\endoldtabular}


%\NoAutoSpacing %french autospacing after ":"

\def\graphpath{}

\newcommand{\changegraphpath}[1]{\def\graphpath{#1}}


\newcommand{\vgraphpage}[2][.84\textheight]{%
%	\begin{center}%
		\includegraphics[height=#1]{\graphpath #2}%
%	\end{center}%
}

\newcommand{\hgraphpage}[2][\textwidth]{%
%	\begin{center}%
		\includegraphics[width=#1]{\graphpath #2}%
%	\end{center}%
}

\newcommand{\graphpage}[2][]{%
	\includegraphics[#1]{\graphpath #2}%
}

\bibliographystyle{apalike}

\newcommand{\insertbibliography}[2]{
	\appendix
	\section*{Bibliography}
	\nocite{#2}
%	\makeatletter % to change template
%	\setbeamertemplate{headline}[default] % not mandatory, but I though it was better to set it blank
%	\def\beamer@entrycode{\vspace*{-\headheight}} % here is the part we are interested in :)
%	\makeatother
	\begin{multicols*}{2}[\frametitle{\insertsection} \usebeamertemplate*{frametitle}]%\usebeamertemplate*{frametitle}\frametitle{Références}
		\tiny
		\bibliography{#1}
	\end{multicols*}
}

\definecolor{my-grey}{RGB}{233, 233, 233}

\newcommand{\insertlicence}{
	\begin{frame}[plain]
	\frametitle{License: CC-BY 4.0}

	\begin{tcolorbox}[colback=cyan,
		colframe=cyan,  
		arc=0pt,outer arc=0pt,
		valign=top, 
		halign=center,
		width=\textwidth]
		
		\includegraphics[width=.5cm]{../../img/licence/cc_icon_white_x2.png}
		\includegraphics[width=.5cm]{../../img/licence/attribution_icon_white_x2.png}
		
		\color{white}
		\bfseries Attribution 4.0 International (CC BY 4.0) \\
		\tiny \url{https://creativecommons.org/licenses/by/4.0/deed.en}
		
	\end{tcolorbox}\vspace{-.5cm}
	\begin{tcolorbox}[colback=my-grey,
		colframe=my-grey,  
		center, arc=0pt,outer arc=0pt,
		valign=top, 
		halign=left,
		width=\textwidth]
		
		\tiny
		
		\begin{center}
			\bfseries\Large
			You are free to:
		\end{center}
		
		\begin{minipage}{0.83\textwidth}
			\begin{itemize}
				\item[] \textbf{Share} — copy and redistribute the material in any medium or format
				\item[] \textbf{Adapt} — remix, transform, and build upon the material
				for any purpose, even commercially.
			\end{itemize}
		
			\vspace{12pt}
		
			\textit{The licensor cannot revoke these freedoms as long as you follow the license terms.}
		\end{minipage}
		\begin{minipage}{0.15\textwidth}
			\includegraphics[width=\textwidth]{../../img/licence/FreeCulturalWorks_seal_x2.jpg}
		\end{minipage}
	
		
		\begin{center}
			\bfseries\Large
			Under the following terms:
		\end{center}
		
		\begin{itemize}
			\item[] \textbf{Attribution} — You must give appropriate credit, provide a link to the license, and indicate if changes were made. You may do so in any reasonable manner, but not in any way that suggests the licensor endorses you or your use.
			\item[] \textbf{No additional restrictions} — You may not apply legal terms or technological measures that legally restrict others from doing anything the license permits.
		\end{itemize}
		
	\end{tcolorbox}
	
%	\begin{center}
%		\bfseries Attribution 4.0 International (CC BY 4.0)
%		\url{https://creativecommons.org/licenses/by/4.0/deed.fr}
%	\end{center}

%	\tiny
%
%	Vous êtes autorisé à : 
%	\begin{itemize}
%		\item \textbf{Partager} — copier, distribuer et communiquer le matériel par tous moyens et sous tous formats
%		\item \textbf{Adapter} — remixer, transformer et créer à partir du matériel
%	\end{itemize}
%	
%	Selon les conditions suivantes : 
%	\begin{itemize}
%		\item \textbf{Attribution} — Vous devez créditer l'Œuvre, intégrer un lien vers la licence et indiquer si des modifications ont été effectuées à l'Oeuvre. Vous devez indiquer ces informations par tous les moyens raisonnables, sans toutefois suggérer que l'Offrant vous soutient ou soutient la façon dont vous avez utilisé son Oeuvre.
%		\item \textbf{Pas d'Utilisation Commerciale} — Vous n'êtes pas autorisé à faire un usage commercial de cette Oeuvre, tout ou partie du matériel la composant. 
%		\item \textbf{Pas de restrictions complémentaires} — Vous n'êtes pas autorisé à appliquer des conditions légales ou des mesures techniques qui restreindraient légalement autrui à utiliser l'Oeuvre dans les conditions décrites par la licence.
%	\end{itemize}

	\end{frame}
}

\settowidth{\leftmargini}{\usebeamertemplate{itemize item}}
\addtolength{\leftmargini}{\labelsep}

\AtBeginDocument{
	\newcolumntype{L}[2]{>{\vbox to #2\bgroup\vfill\flushleft}p{#1}<{\egroup}} 
	
	\begin{frame}[plain]
		\maketitle
	\end{frame}

	\insertlicence
}


% needs etoolbox; to break links after -
\appto\UrlBreaks{\do\-}


\makeatletter
\newcommand{\xRightarrow}[2][]{\ext@arrow 0359\Rightarrowfill@{#1}{#2}}
\makeatother


\usefonttheme{structurebold}
%\usefonttheme{professionalfonts}


% figure caption
\setlength\abovecaptionskip{2pt}

% equations
\AtBeginDocument{
\abovedisplayskip=0pt
\abovedisplayshortskip=0pt
\belowdisplayskip=0pt
\belowdisplayshortskip=0pt
}


\subtitle[08- Sentence semantics]{Chapter 08\\Sentence semantics} 

\changegraphpath{../img/sent-sem/}

\begin{document}
	
\begin{frame}
	\frametitle{\inserttitle}
	\framesubtitle{\insertshortsubtitle: Introduction}

	\begin{exampleblock}{Example of sentences}
		\begin{center}
			\Large\bfseries
			My cat caught a mouse with its claws.
			
			My cat caught a mouse with pleasure.
			
			My cat caught a mouse with another cat.
		\end{center}
	\end{exampleblock}
	
	\begin{itemize}
		\item What instrument did the cat use to catch the mouse? How? With who?
		\item The last sentence means: 
		\begin{itemize}
			\item My cat caught two animals: a mouse and another mouse.
			\item The two cats caught a mouse.
		\end{itemize}
	\end{itemize}

\end{frame}

\begin{frame}
	\frametitle{\inserttitle}
	\framesubtitle{\insertshortsubtitle: Some humor}
	
	\begin{center}
		\vgraphpage{humor/humor-ambiguity.jpeg}
	\end{center}
	
\end{frame}

\begin{frame}
	\frametitle{\inserttitle}
	\framesubtitle{\insertshortsubtitle: Plan}

	\begin{multicols}{2}
	%	\small
	\tableofcontents
	\end{multicols}

\end{frame}

%===================================================================================
\section{Semantic roles}
%===================================================================================

\begin{frame}
	\frametitle{\insertshortsubtitle}
	\framesubtitle{\insertsection}
	
\end{frame}

\subsection{Thematic roles}

\begin{frame}
	\frametitle{\insertshortsubtitle: \insertsection}
	\framesubtitle{\insertsubsection}
	
	\begin{table}
		\tiny\bfseries
		\begin{tblr}{
				colspec = {p{.13\textwidth}p{.42\textwidth}p{.35\textwidth}},
				row{odd} = {lightblue},
				row{even} = {lightyellow},
				row{1} = {darkblue},
				rowsep=2pt,
			}
			\textcolor{white}{Role} & \textcolor{white}{Description} & \textcolor{white}{Example}\\
			
			AGENT &
			The volitional causer of an event &
			\expword{\underline{Karim} broke the window with a rock.}\\
			
			EXPERIENCER & 
			The experiencer of an event & 
			\expword{\underline{Karim} has a headache.}\\
			
			FORCE &
			The non-volitional causer of the event &
			\expword{\underline{The wind} blows the debris.}\\
			
			THEME &
			The participant most directly affected by an event &
			\expword{Karim broke \underline{the window} with a rock.}\\
			
			RESULT &
			The end product of an event &
			\expword{The city built \underline{a baseball field}.}\\
			
			CONTENT &
			The proposition or content of a propositional event &
			\expword{Sara asked\newline \underline{``Did you meet Leila in a supermarket?"}}\\
			
			INSTRUMENT &
			An instrument used in an event &
			\expword{\underline{a rock} broke the window.}\\
			
			BENEFICIARY &
			The beneficiary of an event &
			\expword{Sara makes hotel reservations for \underline{her boss}.}\\
			
			SOURCE &
			The origin of the object of a transfer event &
			\expword{I came from \underline{Jijel}.}\\
			
			GOAL &
			The destination of an object of a transfer event &
			\expword{I went to \underline{Algiers}.}\\
			
			LOCATIVE & 
			The specification of the place where the action or event designated by the predicate takes place &
			\expword{I live in \underline{Jijel}.}\\
		\end{tblr}
		\caption{Some thematic roles, adapted from \cite{2019-jurafsky-martin}}
	\end{table}
	
\end{frame}

\subsection{FrameNet}

\begin{frame}
	\frametitle{\insertshortsubtitle: \insertsection}
	\framesubtitle{\insertsubsection}

	\begin{minipage}{.68\textwidth}
		\begin{itemize}
			\item {\scriptsize \url{https://framenet.icsi.berkeley.edu/fndrupal/}}
			\item {\scriptsize \url{https://www.nltk.org/howto/framenet.html}}
			\item Resource used to represent meaning
			\item based on``\keyword{Frame semantics}" theory of \keyword{Fillmore}
			\item annotated manually
			\item Many semantic roles
		\end{itemize}
	\end{minipage}
	\begin{minipage}{.3\textwidth}
		\hgraphpage{frameNet-logo.jpg}
	\end{minipage}
	
	\begin{exampleblock}{Example of sentences having the same frame}
		The price of petrol increased.
		
		The price of petrol rose.
		
		There has been a rise in the price of petrol.
	\end{exampleblock}

\end{frame}

\begin{frame}
	\frametitle{\insertshortsubtitle: \insertsection}
	\framesubtitle{\insertsubsection: Structure}
	
	\begin{itemize}
		\item \optword{Frame}: schematic relation of a situation.
		\begin{itemize}
			\item name, definition, semantic type, frame elements, lexical units, examples and relations with other frames
		\end{itemize}
	
		\item \optword{Frame Element (FE)}: a semantic role specific to a frame. 
		\begin{itemize}
			\item It describes a participant or a situation of the frame.
			\item Two types of FE: important (\keyword{Core}) and secondary (\keyword{Non-Core}).
		\end{itemize}
	
		\item \optword{Lexical Units (LU)}: lemmas with their PoS.
		\begin{itemize}
			\item A lexical unit triggers the frame when encountered.
		\end{itemize}
	
		\item \optword{Relations}: inter-frames relations.
		\begin{itemize}
			\item Example, \expword{inheritance}.
		\end{itemize}
	\end{itemize}
	
\end{frame}

\begin{frame}
	\frametitle{\insertshortsubtitle: \insertsection}
	\framesubtitle{\insertsubsection: Example of lexical units}

	\vspace{-12pt}
	\begin{table}
		\scriptsize\bfseries
		\begin{tblr}{
				colspec = {p{.15\textwidth}p{.3\textwidth}p{.45\textwidth}},
				row{odd} = {lightblue},
				row{even} = {lightyellow},
				row{1} = {darkblue},
			}
		
			\textcolor{white}{Lexical Unit} & \textcolor{white}{Frame} & \textcolor{white}{Example}\\
	
			break.n & Opportunity & \\	
			break.v & Cause\_harm & \expword{Jolosa broke a rival player's jaw.}\\
			break.v & Compliance & \expword{He broke his promess.}\\
			break.v & Experience\_bodily\_harm & \expword{I broke my arm in the accident.}\\
			break.v & Cause\_to\_fragment & \expword{Michael broke the bottle against his head}\\
			break.v & Render\_nonfunctional & \expword{I guess I broke the doorknob by twisting it too hard.}\\
			break.v & Breaking\_off & \expword{The handle broke off of the pot.}\\
			break.v & Breaking\_apart & \expword{The handle broke off of the pot.}\\
	
		\end{tblr}
		\caption{Sematic frames activated by `break" LU.}
	\end{table}

\end{frame}

\begin{frame}
	\frametitle{\insertshortsubtitle: \insertsection}
	\framesubtitle{\insertsubsection: Example of semantic frame (1)}
	
	
	\begin{table}
		\tiny\bfseries
		\begin{tblr}{
				colspec = {p{.15\textwidth}p{.75\textwidth}},
				row{odd} = {lightblue},
				row{even} = {lightyellow},
				row{1,3,8} = {darkblue},
				rowsep=0pt,
				cell{1,3,8}{1} = {c=2}{c},
			}
		
			\textcolor{white}{Cause\_to\_fragment} & \\
			
			Definition & An \textcolor{red}{Agent} suddenly and often violently separates the \textcolor{red}{Whole\_patient} into two or more smaller \textcolor{red}{Pieces}, resulting in the \textcolor{red}{Whole\_patient} no longer existing as such. Several lexical items are marked with the semantic type Negative, which indicates that the fragmentation is necessarily judged as injurious to the original \textcolor{red}{Whole\_patient}. Compare this frame with Damaging, Render\_non-functional, and Removing. \\	
			
%			\rowcolor{darkblue}
			\textcolor{white}{FEs (Core)} & \\
			
			Agent [Agt] \newline \textcolor{blue}{Semantic Type: Sentient} & 
			The conscious entity, generally a person, that performs the intentional action that results in the \textcolor{red}{Whole\_patient} being broken into \textcolor{red}{Pieces}. \newline \expword{\underline{I and I alone} can SHATTER the gem and break the curse.} \\
			
			Cause [cau] & 
			An event which leads to the fragmentation of the \textcolor{red}{Whole\_patient}. \\
			
			Pieces [Pieces]	& 
			The fragments of the \textcolor{red}{Whole\_patient} that result from the \textcolor{red}{Agent}'s action.
			\newline
			\expword{I SMASHED the toy boat to \underline{flinders}.} \\
			
			Whole\_patient [Pat] & The entity which is destroyed by the \textcolor{red}{Agent} and that ends up broken into \textcolor{red}{Pieces}.
			\newline
			\expword{Shattering someone's confidence is a little different than SHATTERING \underline{a dish}.} \\
			
%			\rowcolor{darkblue}
			\textcolor{white}{FEs (None-Core)} & \\
			
			Degree [Degr] \newline \textcolor{blue}{Semantic Type: Degree} &
			The degree to which the fracturing is completed. 
			\newline
			\expword{I SHATTERED the vase \underline{completely}.} \\
			
			Explanation [Exp] \newline \textcolor{blue}{Semantic Type: State\_of\_affairs} &
			A state of affairs that the Agent is responding to in performing the action. \newline
			\expword{He TORE the treaty UP out of frustration.}
			
		\end{tblr}
		\caption{Example of frame ``Cause\_to\_fragment" (part 1)
			\newline
			{\tiny\url{ https://framenet2.icsi.berkeley.edu/fnReports/data/frameIndex.xml?frame=Cause_to_fragment}}%
		}
	\end{table}
	
\end{frame}

\begin{frame}
	\frametitle{\insertshortsubtitle: \insertsection}
	\framesubtitle{\insertsubsection: Example of semantic frame (2)}
	
	\begin{table}
		\tiny\bfseries
		\begin{tblr}{
				colspec = {p{.15\textwidth}p{.75\textwidth}},
				row{odd} = {lightblue},
				row{even} = {lightyellow},
				row{1,3,7,11} = {darkblue},
				rowsep = 0pt,
				cell{1,3,6,7,11}{1} = {c=2}{c},
			}
			\textcolor{white}{Cause\_to\_fragment} & \\
			
			 & \\	
			 
			
%			\rowcolor{darkblue}
			\textcolor{white}{FEs (None-Core)} & \\
			
			Explanation [Exp] \newline \textcolor{blue}{Semantic Type: State\_of\_affairs} &	
			A state of affairs that the \textcolor{red}{Agent} is responding to in performing the action.
			\newline
			\expword{He TORE the treaty UP \underline{out of frustration}.} \\
			
			Instrument [Ins] \newline \textcolor{blue}{Semantic Type: Physical\_entity} &
			An entity directed by the  \textcolor{red}{Agent} that interacts with a \textcolor{red}{Whole\_patient} to accomplish its fracture. \\
			
			
			{\large ...} & \\
			
%			\rowcolor{darkblue}
			\textcolor{white}{Frame-frame Relations} & \\
			
			Inherits from & Transitive\_action \\
			Uses & Destroying \\
			Is Causative of & Breaking\_apart \\
			
%			\rowcolor{darkblue}
			\textcolor{white}{Lexical Units} & \\
			
			& break apart.v, break down.v, break up.v, break.v, chip.v, cleave.v, dissect.v, dissolve.v, fracture.v, fragment.v, rend.v, rip up.v, rip.v, rive.v, shatter.v, shiver.v, shred.v, sliver.v, smash.v, snap.v, splinter.v, split.v, take apart.v, tear up.v, tear.v \\
			
		\end{tblr}
		\caption{Example of frame ``Cause\_to\_fragment" (part 2)
			\newline
			{\tiny\url{ https://framenet2.icsi.berkeley.edu/fnReports/data/frameIndex.xml?frame=Cause_to_fragment}}%
		}
	\end{table}
	
\end{frame}


\begin{frame}
	\frametitle{\insertshortsubtitle: \insertsection}
	\framesubtitle{\insertsubsection: Example of valence patterns}
	
	\vspace{-6pt}
	\begin{table}
		\tiny\bfseries
		\begin{tabular}{|p{.12\textwidth}|p{.12\textwidth}|p{.12\textwidth}|p{.12\textwidth}|p{.12\textwidth}|p{.12\textwidth}|}
			\hline
			\rowcolor{darkblue}
			\textcolor{white}{Number Annotated} & \multicolumn{5}{|l|}{\textcolor{white}{Patterns}}\\
			\hline
			\multicolumn{6}{l}{ }\\
			
			\hline
			\rowcolor{lightyellow}
			1 TOTAL & \textcolor{red}{Agent} & \textcolor{red}{Instrument} & \textcolor{red}{Pieces} & \textcolor{red}{Whole\_patient} & \\
			\hline
			\rowcolor{lightyellow}
			(1) & CNI \newline - - & PP[with] \newline Dep & INI \newline - - & NP \newline Ext & \\
			\hline
			\multicolumn{6}{l}{ }\\
			
			\hline
			\rowcolor{lightblue}
			1 TOTAL & \textcolor{red}{Agent} & \textcolor{red}{Means} & \textcolor{red}{Pieces} & \textcolor{red}{Time} & \textcolor{red}{Whole\_patient} \\
			\hline
			\rowcolor{lightblue}
			(1) & NP \newline Ext & 2nd \newline - - & INI \newline - - & Sinterrog \newline Dep & NP \newline Obj \\
			\hline
			\multicolumn{6}{l}{ }\\
			
			\hline
			\rowcolor{lightyellow}
			4 TOTAL & \textcolor{red}{Agent} & \textcolor{red}{Pieces} & \textcolor{red}{Whole\_patient} & & \\
			\hline
			\rowcolor{lightyellow}
			(4) & NP \newline Ext & INI \newline - - & NP \newline Obj & & \\
			\hline
		\end{tabular}
		\caption{Lexical entry of trigger ``fracture" (verb) of frame ``Cause\_to\_fragment": valence pattern list excerpt.}
	\end{table}
	
\end{frame}

\begin{frame}
	\frametitle{\insertshortsubtitle: \insertsection}
	\framesubtitle{\insertsubsection: Example of lexicographic annotations}
	
	\vskip-5pt
	\begin{figure}
		\begin{tcolorbox}[colback=white, colframe=blue, boxrule=1pt, text width=.91\textwidth]
			\tiny\bfseries
			\vspace{-6pt}
		\begin{itemize}
			\item 429-s20-rcoll-skull
			\begin{enumerate}\tiny
				\item \ [\textsubscript{\color{red}Agent} Former England Under-21 player Keith Benton] FRACTURED\textsuperscript{\color{red}Target} [\textsubscript{\color{red}Whole\_patient} his son Seb 's skull] [\textsubscript{\color{red}Time} when he hit the ball into the crowd during a match in Buckingham]. [\textsubscript{\color{red}Pieces} INI] 
				\item \ [\textsubscript{\color{red}Agent} He] hit a lamp-post and FRACTURED\textsuperscript{\color{red}Target} [\textsubscript{\color{red}Whole\_patient} Mike 's skull]. [\textsubscript{\color{red}Pieces} INI] 
				\item When he found the man [\textsubscript{\color{red}Agent} he] threw the acid into his face and beat him with the hammer , FRACTURING\textsuperscript{\color{red}Target} [\textsubscript{\color{red}Whole\_patient} his skull] and his thumb. [\textsubscript{\color{red}Pieces} INI] 
				\item \ [\textsubscript{\color{red}Agent} A nanny] has been jailed after FRACTURING\textsuperscript{\color{red}Target} [\textsubscript{\color{red}Whole\_patient} the skulls of two new born babies in her care]. [\textsubscript{\color{red}Pieces} INI] 
			\end{enumerate}
			\item 520-s20-np-vping
			\item 620-s20-np-ppother
			\item 660-s20-trans-simple
			\begin{enumerate}\tiny
				\item \ [\textsubscript{\color{red}Agent} Then 17-year-old Lee Diaz, of North End Gardens, Bishop Auckland], attacked a second party-goer, Carl Gent, punching him in the face and FRACTURING\textsuperscript{\color{red}Target} [\textsubscript{\color{red}Whole\_patient} his jaw]. [\textsubscript{\color{red}Pieces} INI] 
			\end{enumerate}
			
%			\item 670-s20-pass-by
			\item 680-s20-pass
			\begin{enumerate}\tiny
				\item \ [\textsubscript{\color{red}Whole\_patient} It] was FRACTURED\textsuperscript{\color{red}Target} [\textsubscript{\color{red}Instrument} with a solvent-cleaned chisel], and the outer orange layer discarded. [\textsubscript{\color{red}Agent} CNI][\textsubscript{\color{red}Pieces} INI] 
			\end{enumerate}
			
%			\item 690-s20-trans-other
%			\item 730-s20-ppwith
%			\item 780-s20-ppother
%			\item 880-s20-intrans-simple
%			\item 890-s20-intrans-adverb
%			\item 900-s20-other
		\end{itemize}\vspace{-6pt}
	\end{tcolorbox}\vspace{-6pt}

		\caption{Lexicographic annotations excerpt of the trigger ``fracture" (verb) of the frame ``Cause\_to\_fragment".}
	\end{figure}
	
\end{frame}

\subsection{PropBank}

\begin{frame}
	\frametitle{\insertshortsubtitle: \insertsection}
	\framesubtitle{\insertsubsection}

	
	%\begin{minipage}{.68\textwidth}
	\begin{itemize}
		\item {\scriptsize \url{https://propbank.github.io/}}
		\item {\scriptsize \url{https://www.nltk.org/howto/propbank.html}}
		\item Annotated corpus based on Predicate-Arguments structure
		\item Less semantic roles 
		\begin{itemize}
			\item \optword{Proto-Agent}
			\begin{itemize}
				\item voluntarily participate in an event or state
				\item cause an event or a state change of another participant
			\end{itemize}
			\item \optword{Proto-Patient}
			\begin{itemize}
				\item experience a change of state
				\item be affected by another participant
			\end{itemize}
		\end{itemize}
	\end{itemize}
	
\end{frame}

\begin{frame}
	\frametitle{\insertshortsubtitle: \insertsection}
	\framesubtitle{\insertsubsection: Structure}
	
	\begin{itemize}
		\item \optword{Roleset id}: Verbs are annotated by their meanings.
		\begin{itemize}
			\item E.g. \expword{know.01: be cognizant of, realize ; know.02: be familiar with, have experienced}		
		\end{itemize}
		
		\item \optword{Roles}: Each verb/meaning has a set of possible arguments.  
		\begin{itemize}
			\item \optword{Arg0}: PROTO-AGENT
			\item \optword{Arg1}: PROTO-PATIENT
			\item \optword{Arg2}: in general; beneficiary, instrument, attribute, or end state.
			\item \optword{Arg3}: in general; start point, beneficiary, instrument, or attribute.
			\item \optword{Arg4}: end point
			\item Arg2 ... Arg5 are not consistent in the corpus
		\end{itemize}
	
		\item \optword{Modifiers}: Marked by \keyword{ArgM}  
		\begin{itemize}
			\item \optword{ArgM-TMP}: When?
			\item \optword{ArgM-LOC}: Where?
			\item \optword{ArgM-MNR}: How? 
			\item ...
		\end{itemize}
		
		\item \optword{Annotated examples} 
	\end{itemize}
	
\end{frame}

\begin{frame}
	\frametitle{\insertshortsubtitle: \insertsection}
	\framesubtitle{\insertsubsection: Example of predicates}
	

	\begin{figure}
	%	\tiny
		\scriptsize
		\begin{tcolorbox}[colback=white, colframe=blue, boxrule=1pt, text width=.91\textwidth]
		\begin{itemize}
			\item \textbf{Roleset id}
			\begin{itemize}\scriptsize
				\item \textbf{know.01}: be cognizant of, realize
			\end{itemize}
			\item \textbf{Roles}
			\begin{itemize}\scriptsize
				\item \textbf{Arg0}: knower
				\item \textbf{Arg1}: fact that is known
				\item \textbf{Arg2}: entity that arg1 is known ABOUT
			\end{itemize}
		
			\item \textbf{Example: know-v: sentential thing known}
			\begin{itemize}\scriptsize
				\item \ [\textsubscript{\color{red}Arg0} The other side] knows [\textsubscript{\color{red}Arg1} that Giuliani has always been prochoice].
			\end{itemize}
		
			\item \textbf{Example: know-v: attributive}
			\begin{itemize}\scriptsize
				\item \ [\textsubscript{\color{red}Arg0} He] did[\textsubscript{\color{red}ArgM-NEG} n't] know [\textsubscript{\color{red}Arg1} (anything)] [\textsubscript{\color{red}Arg2} about most of the cases] [\textsubscript{\color{red}ArgM-TMP} until Wednesday].
			\end{itemize}
		\end{itemize}
		\end{tcolorbox}
			
		\caption{ProBank annotation excerpt of predicate ``know" \url{http://verbs.colorado.edu/propbank/framesets-english-aliases/know.html}}
	\end{figure}
	
\end{frame}


%===================================================================================
\section{Semantic roles labeling (SRL)}
%===================================================================================

\begin{frame}
	\frametitle{\insertshortsubtitle}
	\framesubtitle{\insertsection}

	\begin{figure}
		\hgraphpage{exp-srl_.pdf}
		\caption{Example of ProbBank-based SRL [\url{https://demo.allennlp.org/semantic-role-labeling/}]}
	\end{figure}
	
\end{frame}

\subsection{Features-based}

\begin{frame}
	\frametitle{\insertshortsubtitle: \insertsection}
	\framesubtitle{\insertsubsection}
	
	\begin{itemize}
		\item \optword{Idea}: browse the parse tree and classify the nodes.
		\begin{itemize}
			\item \optword{Input}: features on the node (phrase) and the predicate.
			\item \optword{Output}: classes are those of FrameNet or PropBank + \keyword{None} to mark a node with no role
		\end{itemize} 
		\item The classification process can be optimized
		\begin{enumerate}
			\item \optword{Pruning}: eliminate classification constituents using heuristics.
			\item \optword{Identification}: classify the nodes in \keyword{Argument} or \keyword{None}.
			\item \optword{Classification}: classes are those of FrameNet or PropBank.
		\end{enumerate} 
	\end{itemize}
	
\end{frame}

\begin{frame}
	\frametitle{\insertshortsubtitle: \insertsection}
	\framesubtitle{\insertsubsection: Example of annotated parse tree}
	
	\vspace{-0.2cm}
	\begin{figure}
		\hgraphpage{srl-tree.pdf}
		\caption{Example of a parse tree with a path from a phrase to the main predicate \cite{2019-jurafsky-martin}}
	\end{figure}
	
\end{frame}

\begin{frame}
	\frametitle{\insertshortsubtitle: \insertsection}
	\framesubtitle{\insertsubsection: Some features}

	\begin{itemize}
		\item Sentence's main predicate. E.g. \expword{Using computers \underline{enhances} work.};
		\item Phrase type. E.g. \expword{NP, S, PP};
		\item Phrase head (main) word. E.g. \expword{The extremely good \underline{weather}}; 
		\item Head word's PoS. E.g. \expword{The extremely good weather/\underline{NOUN}};
		\item Path from the current node to the predicate. E.g. \expword{NP\textuparrow S\textdownarrow VP \textdownarrow VPD};
		\item Voice: active or passive;
		\item Relative position to the predicate: before or after;
		\item ...
	\end{itemize}
	
\end{frame}

\subsection{Neural-networks-based}

\begin{frame}
	\frametitle{\insertshortsubtitle: \insertsection}
	\framesubtitle{\insertsubsection}

	\begin{minipage}{.48\textwidth}
		\begin{itemize}
			\item \optword{Input}: Words embeddings + predicate indicator (predicate or not).
			\item \optword{Output}: PropBank senses probabilities.  
			\item Using \keyword{IOB} notation
		\end{itemize}
	\end{minipage}
	\begin{minipage}{.5\textwidth}
		\begin{figure}
			\hgraphpage{srl-lstm.pdf}
			\caption{SRL system using LSTM \cite{2017-he-al}}
		\end{figure}
	\end{minipage}
	
\end{frame}

%===================================================================================
\section{Semantic representation of sentences}
%===================================================================================

\begin{frame}
	\frametitle{\insertshortsubtitle: \insertsection}
	\framesubtitle{\insertsubsection}
	
	\begin{itemize}
		\item There are several ways to express the same meaning.
		\begin{itemize}
			\item \expword{The student has prepared a report.}
			\item \expword{A report has been prepared by the student.}
		\end{itemize}
	
		\item Sentences in natural languages can be ambiguous
		\begin{itemize}
			\item \expword{John floated the boat between the rocks.}
			\item See more: \url{https://plato.stanford.edu/entries/ambiguity/}
		\end{itemize}
	
		\item Multiple tasks can benefit from semantic representation:
		\begin{itemize}
			\item Natural language understanding (NLU)
			\item Question-Answering
			\item Information retrieval (IR)
			\item Machine translation
			\item Automatic text summarization
		\end{itemize}
	\end{itemize}
	
\end{frame}

\subsection{First order logic}

\begin{frame}
	\frametitle{\insertshortsubtitle: \insertsection}
	\framesubtitle{\insertsubsection}
	
	\begin{itemize}
		\item \optword{term}: an object.
		\begin{itemize}
			\item \optword{constant}: a specific object
			%
			\expword{Karim, ESI, Algeria}
			
			\item \optword{function}: return an object based on others
			%
			\expword{LocationOf(ESI)}
			
			\item \optword{variable}: anonymous object 
			%
			\expword{x, y, z}
			
		\end{itemize}
		\item \optword{predicate}: a relation 
		
		\expword{University(ESI), TeacherAt(Karim, ESI)}
		
		\item \optword{connective}: AND ($ \wedge $), OR ($ \vee $), NOT ($ \neg $ ), IMPLY ($\rightarrow$) and EQUIVALENT ($ \Leftrightarrow $)
		
		\expword{TeacherAt(Karim, ESI) $\wedge$ University(ESI)}
		
		\item \optword{quantifier}: EXISTS ($\exists$) and ALL ($\forall$)
		
		\expword{I eat at a restaurant near ESI.}
		
		\expword{$\exists$ x Restaurant(x) $\wedge$ Near(LocationOf(x), LocationOf(ESI)) $\wedge$  Eat(Interlocutor, x)}
	\end{itemize}
	
\end{frame}

\begin{frame}
	\frametitle{\insertshortsubtitle: \insertsection}
	\framesubtitle{\insertsubsection: Formalism}
	
	\begin{figure}
%		\vgraphpage[0.6\textheight]{LPO-gram_.pdf}
		\scriptsize
		\begin{tabular}{rcl}
			\hline\hline
			Formula & \textrightarrow & AtomicFormula \\
			        & \textbar        & Formula Connective Formula \\
			        & \textbar        & Quantifier Variable, ... Formula \\
			        & \textbar        & $\textlnot$ Formula \\
			        & \textbar        & (Formula) \\
			AtomicFormula & \textrightarrow & Predicate (Term, ...) \\
			Term    & \textrightarrow & Function(Term, ...) \\
			        & \textbar        & Constant \\
			        & \textbar        & Variable \\
			Connective & \textrightarrow & $\wedge$ \textbar $\vee$ \textbar $\Rightarrow$ \\
			Quantifier & \textrightarrow & $\forall$ \textbar $\exists$ \\
			Constant & \textrightarrow & ESI \textbar Karim \textbar Algérie ...\\
			Variable & \textrightarrow & x \textbar y \textbar ... \\
			Predicate & \textrightarrow & Ecole \textbar EnseignatA \textbar Utiliser \textbar ... \\
			Function & \textrightarrow & EmplacementDe \textbar ... \\
			\hline\hline
		\end{tabular}
		\caption{A grammar specifying first order logic syntax \cite{2019-jurafsky-martin} (Adapted from \cite{2002-russell-norvig})}
	\end{figure}
	
\end{frame}

\begin{frame}
	\frametitle{\insertshortsubtitle: \insertsection}
	\framesubtitle{\insertsubsection: Some humor}
	
	\begin{center}
		\vgraphpage{humor/humor-FOL.jpeg}
	\end{center}
	
\end{frame}

\subsection{Graphs (AMR)}

\begin{frame}[fragile]
	\frametitle{\insertshortsubtitle: \insertsection}
	\framesubtitle{\insertsubsection}
	
	\begin{itemize}
		\item Abstract Meaning Representation \cite{2013-banarescu-al}
		\item A rooted, labeled, directed and acyclic graph
		\item It uses \keyword{PropBank} concepts
	\end{itemize}

	\begin{exampleblock}{Threes formats: Example ``The boy wants to go"}
		\begin{minipage}{.3\textwidth}
			\optword{Logic format}
			
			\footnotesize
			$ \exists $ w, b, g : 
			
			instance(w, want-01) 
			
			$ \wedge $ instance(g, go-01) 
			
			$ \wedge $ instance(b, boy) 
			
			$ \wedge $ arg0(w, b) 
			
			$ \wedge $ arg1(w, g) 
			
			$ \wedge $ arg0(g, b)
		\end{minipage}
		\begin{minipage}{.35\textwidth}
			\optword{AMR format}
			
			\begin{verbatim}
(w / want-01
    :arg0 (b / boy)
    :arg1 (g / go-01
    :arg0 b))
			\end{verbatim}
			
		\end{minipage}
		\begin{minipage}{.3\textwidth}
			\optword{Graph format}
			
			\hgraphpage{amr-graph-exp.pdf}
		\end{minipage}
	\end{exampleblock}
	
\end{frame}

\begin{frame}[fragile]
	\frametitle{\insertshortsubtitle: \insertsection}
	\framesubtitle{\insertsubsection: Relations}
	
	\begin{itemize}
		\item \optword{frame arguments}: following PropBank
		\begin{itemize}
			\item :arg0, :arg1, :arg2, :arg3, :arg4, :arg5
		\end{itemize}
		\item \optword{General semantic relations}
		\begin{itemize}
			\item :accompanier, :age, :beneficiary, :cause, :compared-to, :concession, :condition, :consist-of, :degree, :destination, :direction, :domain, :duration, :employed-by, :example, :extent, :frequency, :instrument, :li, :location, :manner, :medium, :mod, :mode, :name, :part, :path, :polarity, :poss, :purpose, :source, :subevent, :subset, :time, :topic, :value
		\end{itemize}
		\item \optword{Quantity relations}
		\begin{itemize}
			\item :quant, :unit, :scale
		\end{itemize}
		\item \optword{Date relations}
		\begin{itemize}
			\item :day, :month,
			:year, :weekday, :time, :timezone, :quarter,
			:dayperiod, :season, :year2, :decade, :century,
			:calendar, :era
		\end{itemize}
		\item \optword{List relations}
		\begin{itemize}
			\item :op1, :op2, :op3, :op4, :op5,
			:op6, :op7, :op8, :op9, :op10
		\end{itemize}
	\end{itemize}
	
\end{frame}


\subsection{Semantic parsing}

\begin{frame}
	\frametitle{\insertshortsubtitle: \insertsection}
	\framesubtitle{\insertsubsection: Lambda notation}
	
	\begin{itemize}
		\item \keyword{$ \lambda $-Expression}
		\begin{itemize}
			\item Anonymous function
			\item \keyword{$ \lambda $x.P(x)}
			\item Example: \expword{$ \lambda $x.USE(x, BERT)} to express the clause \expword{x uses BERT}
		\end{itemize}
		
		\item \keyword{$ \lambda $-Reduction}
		\begin{itemize}
			\item Substitute a variable by an expression
			\item \keyword{$ \phi (\psi) $}
			\item Example: \expword{$ \lambda $x.USE(x, BERT)\ KARIM = USE(KARIM, BERT)}
		\end{itemize}
		
		\item For fun: \url{http://davidmarino.nfshost.com/lambdalinguist/index.html}
	\end{itemize}
	
\end{frame}

\begin{frame}
	\frametitle{\insertshortsubtitle: \insertsection}
	\framesubtitle{\insertsubsection: Logical representation generation}
	
	\begin{center}
		\small\bfseries
		\begin{tabular}{llll}
			\hline\hline
			S  & \textrightarrow\ NP VP && VP.sem(NP.sem) \\
			VP & \textrightarrow\ V\textsubscript{t} NP && V\textsubscript{t}.sem(NP.sem)\\
			VP & \textrightarrow\ V\textsubscript{i} && V\textsubscript{i}.sem \\
			V\textsubscript{t}  & \textrightarrow\ uses && $ \lambda $y.$ \lambda $x.USE(x, y) \\
			V\textsubscript{i}  & \textrightarrow\ sleeps && $ \lambda $x.SLEEP(x) \\
			NP  & \textrightarrow\  Karim && KARIM \\
			NP  & \textrightarrow\  BERT && BERT \\
			\hline\hline
		\end{tabular}
		
		\hgraphpage[0.7\textwidth]{sem-tree.pdf}
	\end{center}
\end{frame}

\begin{frame}
	\frametitle{\insertshortsubtitle: \insertsection}
	\framesubtitle{\insertsubsection: Quantifiers}
	
	\begin{center}
		\small\bfseries
		\begin{tabular}{llll}
			\hline\hline
			S  & \textrightarrow\ NP VP && NP.sem(VP.sem) \\
			DET & \textrightarrow\ each && $\lambda P.\lambda Q.\forall x (P(x) \Rightarrow Q(x))$ \\
			
			VP & \textrightarrow\ V\textsubscript{t} NP && V\textsubscript{t}.sem(NP.sem) \\
			V\textsubscript{t}  & \textrightarrow\ uses && $\lambda P.\lambda x.P(\lambda y.USE(x, y))$ \\
			
			VP & \textrightarrow\ V\textsubscript{i} && V\textsubscript{i}.sem \\
			V\textsubscript{i}  & \textrightarrow\ sleeps && $ \lambda $x.SLEEP(x) \\
			
			NP & \textrightarrow\ DET NN && DET.sem(NN.sem) \\
			NN  & \textrightarrow\  student && STUDENT \\
			
			NP & \textrightarrow\ NNP && $\lambda P.P(NNP.sem)$ \\
			NNP  & \textrightarrow\  Karim && KARIM \\
			
			DET & \textrightarrow\ a && $\lambda P.\lambda Q.\exists x\ P(x) \wedge Q(x)$ \\
			NNP  & \textrightarrow\  BERT && BERT \\
			\hline\hline
		\end{tabular}
	\end{center}
	
\end{frame}

\begin{frame}
	\frametitle{\insertshortsubtitle: \insertsection}
	\framesubtitle{\insertsubsection: Quantifiers (Example)}
	
	\begin{center}
		\hgraphpage[0.8\textwidth]{sem-qtree.pdf}
		\small
		\begin{align*}
			VP.sem & = V_t.sem(NP.sem) \\
			& = \textcolor{red}{\lambda P}.\lambda x.\textcolor{red}{P}(\lambda y.USE(x, y))(\textcolor{blue}{\lambda P.P(BERT)}) \\
			& = \lambda x.\textcolor{red}{\lambda P}.\textcolor{red}{P}(BERT)(\textcolor{blue}{\lambda y.USE(x, y)}) \\
			& = \lambda x.\textcolor{red}{\lambda y}.USE(x, \textcolor{red}{y})(\textcolor{blue}{BERT}) \\
			& = \lambda x.USE(x, BERT) \\
		\end{align*}
	\end{center}
	
\end{frame}

\begin{frame}[fragile]
	\frametitle{\insertshortsubtitle: \insertsection}
	\framesubtitle{\insertsubsection: NLTK example}
	
	\begin{exampleblock}{Example pf semantic parsing using NLTK}
		\optword{Code:}
		{\scriptsize
			\begin{lstlisting}[language=Python]
from nltk import load_parser
parser = load_parser('grammars/book_grammars/simple-sem.fcfg', trace=0)
sentence = 'Angus gives a bone to every dog'
tokens = sentence.split()
for tree in parser.parse(tokens):
print(tree.label()['SEM'])
			\end{lstlisting}
		}
		
		\optword{Result:}
		{\scriptsize\bfseries
			\begin{lstlisting}
all z2.(dog(z2) -> exists z1.(bone(z1) & give(angus,z1,z2)))
			\end{lstlisting}
		}
		
	\end{exampleblock}
	
	\begin{itemize}
		\item See: \url{https://www.nltk.org/book/ch10.html}
	\end{itemize}
	
\end{frame}

\begin{frame}
	\frametitle{\insertshortsubtitle: \insertsection}
	\framesubtitle{\insertsubsection: Some humor}
	
	\begin{center}
		\vgraphpage{humor/humor-sens.jpg}
	\end{center}
	
\end{frame}

\insertbibliography{NLP08}{*}

\end{document}

